\section{Executable Evidence Contract}
A campaign declares seed, condition, law, fault, horizon $H$, frame duration and resource unit. Adapters normalize native traces into those identities, source/map/calibration bindings, input/mask digests, completed observations, cumulative resource and first failure. Hashes bind bytes, not execution. Common checks test matrix membership, pairing, clocks, exposure, resources and outcome labels; native replay separately regenerates inputs and plant transitions. Shared variates need not produce identical feedback-dependent events. Warning prefixes exclude future labels.

Native timing checks require burst onset 250\,ms after the trigger, 450\,ms duration, and an active-at-$H$ flag consistent with that interval. In 72 separate historical records, an audit found 17 delays of 260\,ms, four completed bursts of 440\,ms and six incorrect horizon flags. These overlapping categories are not pooled with current cohorts. Ignoring them would misstate fast-motion duration or actuation at the endpoint. Replay therefore checks realized events, rather than accepting nominal parameters. Core-schema findings retain row indices; arbitrary native semantics remain adapter obligations.

\reuseRef\ adapts the same checker to a delayed tank: $\dot h=I-u-0.04h$, $h(0)=0.5$, safe interval $[0.2,0.8]$. Eight seeds, two inflows, two laws and three sensor-deletion levels (0/25/50\%) give 96 examples. Deleted samples hold their last value. Independent analytic replay checks sensor accounting and 150-ms delayed actuation on a 100-ms grid. All examples remain safe: they demonstrate reuse, not neural fault tolerance.

The original challenge set detects eleven common-checker and three native-only cases per adapter, not general coverage. A synthetic zero-neural-readout logger passes input/plant replay without solving neural equations: consistent traces cannot authenticate execution, unlogged state or safety.
